% Lösungen Einheit 1 von 2 – Umkehrbarkeit, Bereiche und Term (zusammenbau v0.1)
\einheitenkopf{Einheit 1 von 2 · Umkehrbarkeit, Bereiche und Term}

\erg{8}{a) umkehrbar – f steigt überall, jede waagerechte Gerade trifft höchstens einmal \quad b) $g(x) = \frac{x + 1}{2}$; $g(7) = 4$ \quad c) Definitionsbereich $]3;\, \infty[$ (Wertebereich von f, die 3 wird nicht angenommen), Wertebereich ℝ \quad d) $]-4;\, \infty[$ = Wertebereich von f. $y = 4e^{-\frac{1}{2}x} - 4 \Leftrightarrow e^{-\frac{1}{2}x} = \frac{y}{4} + 1 \Leftrightarrow x = -2\,\mathrm{ln}\left(\frac{y}{4} + 1\right)$, dann x und y tauschen \quad e) Hochpunkt bei $x = 2$ mit $f(2) = \frac{6}{e} \approx 2{,}21$: eine waagerechte Gerade knapp unter dem Hochpunkt trifft den Graphen von f zweimal. Auf $[2;\, \infty[$ fällt f streng, h ist umkehrbar. Umkehrfunktion: Definitionsbereich $]0;\, \frac{6}{e}]$ (0 wird nicht angenommen), Wertebereich $[2;\, \infty[$}
\erg{9}{a) $h = 0{,}75r^2$ nach $r \ge 0$ auflösen: $r = \sqrt{\frac{4h}{3}} = \frac{2}{3}\sqrt{3h}$; für $h = 3$: $r = 2$ cm}
\erg{10}{a) Fehler: der Definitionsbereich von g ist der Wertebereich von f, nicht dessen Definitionsbereich. Richtig: $]-2;\, \infty[$ \quad b) Links und rechts vom Hochpunkt nimmt die Funktion Werte knapp unter dem Hochwert je einmal an; ein solcher Wert gehört zu zwei Stellen, das Rückwärtslesen ist nicht eindeutig. \quad c) $s(P) = \frac{P - 3{,}5}{2{,}1}$; für 24,50 €: $s = 10$ km \quad d) Spalten tauschen: $g(1) = 0$, $g(3) = 1$, $g(5) = 2$, $g(7) = 3$}

10d)

\wertetabelle[0,1,2,3]{x}{g(x)}{1,3,5,7}

